%% ============================================================
%%  1.  INTRODUCTION
%% ============================================================
\section{Introduction}
\label{sec:intro}

Multi-agent LLM debate has emerged as a promising approach to collective
reasoning, in which several language model instances discuss a question and
aim to converge on a correct answer~\cite{du2023improving,liang2023encouraging,
chan2023chateval}.
The underlying hope is that disagreement drives critical examination:
agents challenge each other's reasoning, surface overlooked evidence, and
produce more robust conclusions than any individual agent could.

\paragraph{The majority bias problem.}
In practice, however, convergence often reflects social pressure rather than
epistemic progress.
When a majority of agents align on one answer, minority agents frequently
abandon correct hypotheses not because they have been logically defeated,
but because peer pressure outweighs the strength of their own reasoning.
We call this failure mode \emph{premature elimination}: a correct hypothesis
is present in the debate and then discarded before receiving adequate
evaluation.
Across 200 MMLU-Pro questions under a standard debate protocol, we find that
8.2\% of proposed correct hypotheses are discarded before the final round—
a non-trivial loss rate that directly reduces accuracy.

\paragraph{The observation–control gap.}
Recent work has developed principled process-level diagnostics for detecting
such failures~\cite{diagnostic2026}.
Metrics such as influence asymmetry, balance, and engagement can identify
when a debate is dominated by one agent, converging too rapidly, or
exhibiting dogmatic non-engagement.
However, this prior work \emph{observes} pathological dynamics after the
fact; it does not \emph{correct} them.
Closing this observe–control loop is the central contribution of this paper.

\paragraph{Adaptive Influence Balancing.}
We propose \textbf{Adaptive Influence Balancing} (AIB), a framework that
monitors live process diagnostics and selectively intervenes to preserve
minority agents whose reasoning shows evidence of correctness.
AIB instantiates three complementary mechanisms—loser resistance, static
minority protection, and diagnostic-triggered minority protection—and
evaluates them alongside a learned predictor that replaces hand-tuned
thresholds with a calibrated estimate of minority correctness.

\paragraph{Key results.}
Across 200 MMLU-Pro questions with Qwen2.5-14B-Instruct (3 agents, 5 rounds):
\begin{itemize}[noitemsep,topsep=2pt]
  \item \textbf{Diagnostic-triggered minority protection} (Exp3b) achieves
        60.0\% accuracy vs.\ 54.5\% baseline (+5.5~pp, $p\!<\!0.01$),
        reduces PER from 8.2\% to 3.2\% (−61.3\%), and improves consensus
        reliability to 60.6\%.
  \item \textbf{Learned predictor} (Exp4) replaces fixed thresholds with a
        logistic predictor trained on process diagnostics and reasoning
        quality; top features are minority reasoning quality ($\beta=0.245$)
        and support trajectory ($\beta=0.179$).
  \item \textbf{Not all interventions help}: adversarial loser resistance
        (Exp2a) \emph{worsens} PER by 9.7\%, and categorical minority
        protection (Exp3a) paradoxically increases influence asymmetry.
        Quality-gated, diagnostics-informed triggers are essential.
\end{itemize}

\paragraph{Contributions.}
\begin{enumerate}[noitemsep,topsep=2pt]
  \item The \textbf{AIB framework}: three intervention mechanisms for
        protecting minority hypotheses, ranging from support-triggered
        heuristics to process-diagnostic control signals.
  \item \textbf{Premature Elimination Rate (PER)}: a new metric that
        directly measures the failure mode targeted by AIB and provides a
        more diagnostic signal than accuracy alone.
  \item \textbf{Empirical evidence}: +5.5~pp accuracy and −61\% PER on
        MMLU-Pro, with evidence that trigger selectivity determines whether
        interventions help or hurt.
  \item \textbf{Learned predictor}: replaces hard thresholds with a
        calibrated model showing that reasoning quality and support
        dynamics are the dominant predictors of minority correctness.
\end{enumerate}
